\section{Results Tracking and Internal Notes}

This section replaces the traditional ``Results'' and ``Discussion'' split with a practical development log: what has been generated, what remains pending, and which methodological questions are still open.

\subsection{Generated artifacts status}

The current pipeline emits quality-control and analysis outputs at each stage: ingestion/filtering diagnostics, ICA component summaries and evaluation tables, segmentation/baseline/artifact-rejection statistics, FFT/CSD/wPLI visualizations and matrices, and surrogate-based inference tables. Keeping these outputs phase-scoped is essential for rapid debugging because it allows validation of local assumptions before propagating changes to downstream stages.
\sidenote{%
\footnotesize
\textbf{Minimum per run}\\
1) phase logs\\
2) QC figures\\
3) summary tables\\
4) config snapshot
}

\subsection{Pending tasks}

The main pending items are operational rather than editorial: verify EO/EC completeness and labeling consistency for all participants, lock the final surrogate count and rerun inference under the frozen configuration, and consolidate band-wise outputs into reproducible internal comparison tables. These tasks should be tracked with run timestamps and configuration hashes to avoid ambiguity between exploratory and reference executions.
\sidenote{%
\footnotesize
\textbf{Recommended tag}\\
\texttt{run\_<date>\_<branch>\_<cfg>}\\
Include commit hash and key params.
}

\subsection{Working hypotheses and open questions}

Current hypotheses remain centered on expected EO/EC modulation in alpha-related connectivity, improved specificity after CSD+wPLI relative to non-CSD alternatives, and stabilization of significant-edge density when surrogate thresholds are applied. Open technical questions include sensitivity to segment length/overlap choices, potential over-cleaning in low-SNR subjects during ICA artifact rejection, and runtime-versus-precision trade-offs when increasing the number of surrogates.
